% !TeX program = lualatex
\documentclass[12pt,aspectratio=169]{beamer}
\usepackage[utf8]{inputenc}
\usepackage[T1]{fontenc}
\usepackage{amsmath,amsfonts,amssymb}
\usepackage{graphicx}
\usepackage{xcolor}
\usepackage{hyperref}
\usepackage{anyfontsize}
\usepackage{lmodern}

% ====== EXTREME FONT REDUCTION ======
\renewcommand{\tiny}{\fontsize{4}{5}\selectfont}
\renewcommand{\scriptsize}{\fontsize{5}{6}\selectfont}
\renewcommand{\footnotesize}{\fontsize{6}{7}\selectfont}
\renewcommand{\small}{\fontsize{7}{8}\selectfont}
\renewcommand{\normalsize}{\fontsize{8}{9}\selectfont}
\renewcommand{\large}{\fontsize{9}{10}\selectfont}
\renewcommand{\Large}{\fontsize{10}{11}\selectfont}
\renewcommand{\LARGE}{\fontsize{11}{13}\selectfont}
\renewcommand{\huge}{\fontsize{12}{14}\selectfont}
\renewcommand{\Huge}{\fontsize{14}{17}\selectfont}

% ====== COLOR THEME ======
\definecolor{redcorrect}{RGB}{200,30,30}
\definecolor{bluecorrect}{RGB}{30,80,200}
\definecolor{greencheck}{RGB}{0,150,50}
\definecolor{lightgray}{RGB}{245,245,245}
\definecolor{darkgray}{RGB}{50,50,50}
\definecolor{softyellow}{RGB}{255,248,200}
\definecolor{steppurple}{RGB}{150,50,200}
\definecolor{deepblue}{RGB}{20,60,150}
\definecolor{darkred}{RGB}{150,20,20}
\definecolor{orangewarning}{RGB}{255,140,0}

\setbeamercolor{title}{fg=white,bg=redcorrect}
\setbeamercolor{frametitle}{fg=white,bg=redcorrect}
\setbeamercolor{block title}{fg=white,bg=bluecorrect}
\setbeamercolor{block body}{fg=darkgray,bg=lightgray}
\setbeamercolor{normal text}{fg=darkgray}
\usecolortheme{default}

% ====== CUSTOM COMMANDS ======
\newcommand{\correct}[1]{\textcolor{greencheck}{\textbf{\checkmark}} #1}
\newcommand{\keypoint}[1]{\textcolor{deepblue}{\textbf{#1}}}
\newcommand{\hardconcept}[1]{\textcolor{darkred}{\textbf{#1}}}
\newcommand{\tipbox}[1]{\fcolorbox{bluecorrect}{softyellow}{\parbox{0.88\textwidth}{\centering #1}}}
\newcommand{\stepbox}[1]{%
	\begin{center}
		\fcolorbox{darkgray}{white}{%
			\parbox{0.90\textwidth}{\vspace{0.05cm} #1 \vspace{0.05cm}}%
		}
	\end{center}
}
\newcommand{\wrongbox}[1]{\fcolorbox{redcorrect}{red!10}{\parbox{0.88\textwidth}{\centering \textcolor{redcorrect}{\textbf{\times}} #1}}}
\newcommand{\highlight}[1]{\textcolor{redcorrect}{\textbf{#1}}}
\newcommand{\bluetext}[1]{\textcolor{bluecorrect}{\textbf{#1}}}
\newcommand{\stagebox}[1]{\fcolorbox{steppurple}{white!5}{\parbox{0.88\textwidth}{\centering #1}}}

\begin{document}
	
	% ================================================================
	% SLIDE 1: TITLE
	% ================================================================
	\begin{frame}
		\centering
		\vspace{0.8cm}
		{\fontsize{16}{19}\selectfont \textbf{Understanding the Risks and Methods of Financial Crime}}\\[0.15cm]
		{\fontsize{12}{14}\selectfont Domain A — Complete Q\&A Review}\\[0.3cm]
		{\fontsize{9}{11}\selectfont Understanding the \textit{Regulatory Principles} Behind Each Answer}\\[0.3cm]
		{\fontsize{8}{10}\selectfont Compliance Training Department}
		\vspace{0.2cm}
		\rule{4cm}{0.5pt}
	\end{frame}
	
	% ================================================================
	% SLIDE 2: AGENDA
	% ================================================================
	\begin{frame}
		\frametitle{\fontsize{11}{13}\selectfont Agenda — Domain A Topics}
		\scriptsize
		\begin{enumerate}
			\item Definitions of AML, CFT, Sanctions, Fraud, ABC, Tax Evasion (Q1-Q6)
			\item Financial Systems and Consequences of ML (Q7-Q10)
			\item Predicate Crimes (Q11-Q14)
			\item How ML, TF, Sanctions, Bribery, Tax Evasion, Fraud Manifest (Q15-Q20)
			\item Banking Sector Product Risks (Q21-Q24)
			\item Banking High-Risk Sectors (Q25-Q28)
			\item PEP and High-Risk Customer Risks (Q29-Q32)
			\item Other Financial Sector Product Risks (Q33-Q36)
			\item MSB, PSP, and E-Commerce Risks (Q37-Q40)
			\item Insurance Product Risks (Q41-Q44)
			\item VASP, Cryptoassets, and Related Products (Q45-Q48)
			\item Gaming and Gambling Risks (Q49-Q52)
			\item Real Estate Risks (Q53-Q56)
			\item Gatekeeper Risks (Q57-Q60)
			\item Trust and Company Service Provider Risks (Q61-Q64)
			\item Accountant Risks (Q65-Q68)
			\item Other High-Risk Sectors and Structures (Q69-Q72)
		\end{enumerate}
	\end{frame}
	
	% ================================================================
	% SLIDE 3-8: DEFINITIONS
	% ================================================================
	\begin{frame}
		\frametitle{\fontsize{11}{13}\selectfont Q1: Definition of Money Laundering}
		\begin{block}{\fontsize{8}{10}\selectfont Topic: Definitions of AML, CFT, Sanctions, Fraud, ABC, Tax Evasion}
			\scriptsize
			Which of the following best defines money laundering?
		\end{block}
		\vspace{0.05cm}
		\stepbox{\scriptsize
			\keypoint{Correct Answer:} The process of making illicit funds appear legitimate by concealing their true origin through a series of transactions.
		}
		\vspace{0.05cm}
		\scriptsize
		\hardconcept{Core Concept:} Money laundering has three stages:
		\par \textbf{Placement}: Introducing illicit funds into the financial system.
		\par \textbf{Layering}: Moving funds through complex transactions to obscure origin.
		\par \textbf{Integration}: Returning "cleaned" funds to the legitimate economy.
		\vspace{0.03cm}
		\wrongbox{\scriptsize \textbf{Common Trap:} "Money laundering is simply hiding money" — INCORRECT. It is a \textbf{process} with three distinct stages.}
		\vspace{0.03cm}
		\tipbox{\scriptsize \textbf{Key Insight:} The goal of money laundering is to make illicit funds \textbf{appear legitimate} so they can be used without raising suspicion.}
	\end{frame}
	
	\begin{frame}
		\frametitle{\fontsize{11}{13}\selectfont Q2: Definition of Terrorist Financing}
		\begin{block}{\fontsize{8}{10}\selectfont Topic: Definitions — CFT}
			\scriptsize
			Which of the following best defines terrorist financing?
		\end{block}
		\vspace{0.05cm}
		\stepbox{\scriptsize
			\keypoint{Correct Answer:} The provision or collection of funds, by any means, directly or indirectly, with the intention that they should be used to carry out terrorist acts.
		}
		\vspace{0.05cm}
		\scriptsize
		\hardconcept{Core Concept:} Key differences from ML:
		\par \textbf{Source}: TF funds may be legitimate (charity donations).
		\par \textbf{Amount}: TF can involve small amounts.
		\par \textbf{Intent}: The critical element is \textbf{intent} to support terrorism.
		\par \textbf{Destination}: TF is about \textbf{where} the money is going.
		\vspace{0.03cm}
		\wrongbox{\scriptsize \textbf{Common Trap:} "Terrorist financing always involves illegal money" — INCORRECT. TF funds can come from \textbf{legitimate sources} but are used for illegal purposes.}
		\vspace{0.03cm}
		\tipbox{\scriptsize \textbf{Key Insight:} TF is about \textbf{purpose}, not source. Even legitimate funds can be terrorist financing if intended for terrorism.}
	\end{frame}
	
	\begin{frame}
		\frametitle{\fontsize{11}{13}\selectfont Q3: Definition of Sanctions}
		\begin{block}{\fontsize{8}{10}\selectfont Topic: Definitions — Sanctions}
			\scriptsize
			Which of the following best defines sanctions in the financial crime context?
		\end{block}
		\vspace{0.05cm}
		\stepbox{\scriptsize
			\keypoint{Correct Answer:} Legal restrictions imposed by governments or international bodies against designated countries, individuals, or entities to achieve foreign policy or national security objectives.
		}
		\vspace{0.05cm}
		\scriptsize
		\hardconcept{Core Concept:} Types of sanctions:
		\par \textbf{Country/Regime Sanctions}: Comprehensive restrictions.
		\par \textbf{Targeted Sanctions}: Against specific individuals or entities.
		\par \textbf{Sectoral Sanctions}: Restrictions on specific industries.
		\par \textbf{Asset Freezes}: Freezing assets of designated persons.
		\vspace{0.03cm}
		\wrongbox{\scriptsize \textbf{Common Trap:} "Sanctions are optional recommendations" — INCORRECT. Sanctions are \textbf{legally binding} and violations carry severe penalties.}
		\vspace{0.03cm}
		\tipbox{\scriptsize \textbf{Key Insight:} Sanctions are \textbf{strict liability} — no intent is required for a violation.}
	\end{frame}
	
	\begin{frame}
		\frametitle{\fontsize{11}{13}\selectfont Q4: Definition of Fraud}
		\begin{block}{\fontsize{8}{10}\selectfont Topic: Definitions — Fraud}
			\scriptsize
			Which of the following best defines financial fraud?
		\end{block}
		\vspace{0.05cm}
		\stepbox{\scriptsize
			\keypoint{Correct Answer:} Intentional deception or misrepresentation made for personal gain or to damage another individual, involving financial transactions or assets.
		}
		\vspace{0.05cm}
		\scriptsize
		\hardconcept{Core Concept:} Common fraud types:
		\par \textbf{Identity Theft}: Using stolen personal information.
		\par \textbf{Synthetic Identity Fraud}: Creating fake identities.
		\par \textbf{Investment Fraud}: False investment opportunities.
		\par \textbf{Payment Fraud}: Unauthorized transactions.
		\par \textbf{Insurance Fraud}: False insurance claims.
		\vspace{0.03cm}
		\tipbox{\scriptsize \textbf{Key Insight:} Fraud and money laundering often overlap — fraud generates proceeds that need laundering.}
	\end{frame}
	
	\begin{frame}
		\frametitle{\fontsize{11}{13}\selectfont Q5: Definition of Anti-Bribery and Corruption}
		\begin{block}{\fontsize{8}{10}\selectfont Topic: Definitions — ABC}
			\scriptsize
			Which of the following best defines bribery?
		\end{block}
		\vspace{0.05cm}
		\stepbox{\scriptsize
			\keypoint{Correct Answer:} The offering, giving, receiving, or soliciting of anything of value to influence an action or decision.
		}
		\vspace{0.05cm}
		\scriptsize
		\hardconcept{Core Concept:} Key distinctions:
		\par \textbf{Bribery}: Paying for influence (illegal).
		\par \textbf{Facilitation Payment}: Small payment to speed routine action (allowed under FCPA, prohibited under UK Bribery Act).
		\par \textbf{Hospitality}: Legitimate business entertainment (must be proportionate).
		\par \textbf{Corruption}: Abuse of entrusted power for private gain.
		\vspace{0.03cm}
		\wrongbox{\scriptsize \textbf{Common Trap:} "All payments to officials are bribery" — INCORRECT. Facilitation payments are treated differently under different laws.}
		\vspace{0.03cm}
		\tipbox{\scriptsize \textbf{Key Insight:} The \textbf{UK Bribery Act} is stricter than the FCPA — it prohibits facilitation payments.}
	\end{frame}
	
	\begin{frame}
		\frametitle{\fontsize{11}{13}\selectfont Q6: Definition of Tax Evasion}
		\begin{block}{\fontsize{8}{10}\selectfont Topic: Definitions — Tax Evasion}
			\scriptsize
			Which of the following best defines tax evasion?
		\end{block}
		\vspace{0.05cm}
		\stepbox{\scriptsize
			\keypoint{Correct Answer:} The illegal act of deliberately not paying taxes that are legally owed, typically through concealment or misrepresentation.
		}
		\vspace{0.05cm}
		\scriptsize
		\hardconcept{Core Concept:} Distinction:
		\par \textbf{Tax Evasion}: Illegal — misrepresentation, concealment.
		\par \textbf{Tax Avoidance}: Legal — using legitimate methods to reduce tax liability.
		\vspace{0.03cm}
		\wrongbox{\scriptsize \textbf{Common Trap:} "Tax evasion and tax avoidance are the same" — INCORRECT. Evasion is \textbf{illegal}; avoidance is \textbf{legal}.}
		\vspace{0.03cm}
		\tipbox{\scriptsize \textbf{Key Insight:} Tax evasion is a predicate offense for money laundering in many jurisdictions.}
	\end{frame}
	
	% ================================================================
	% SLIDE 9-12: FINANCIAL SYSTEMS AND CONSEQUENCES
	% ================================================================
	\begin{frame}
		\frametitle{\fontsize{11}{13}\selectfont Q7: Economic Consequences of ML}
		\begin{block}{\fontsize{8}{10}\selectfont Topic: Financial Systems and Consequences of ML}
			\scriptsize
			What is an economic consequence of money laundering?
		\end{block}
		\vspace{0.05cm}
		\stepbox{\scriptsize
			\keypoint{Correct Answer:} Distortion of economic statistics and misallocation of resources, undermining legitimate economic activity.
		}
		\vspace{0.05cm}
		\scriptsize
		\hardconcept{Core Concept:} Economic consequences:
		\par \textbf{Distorted Statistics}: GDP and economic data become unreliable.
		\par \textbf{Misallocation of Resources}: Capital flows to criminal enterprises.
		\par \textbf{Increased Volatility}: Sudden capital movements.
		\par \textbf{Reduced Foreign Investment}: Countries perceived as high-risk.
		\vspace{0.03cm}
		\tipbox{\scriptsize \textbf{Key Insight:} Money laundering \textbf{undermines} legitimate economic activity by distorting markets.}
	\end{frame}
	
	\begin{frame}
		\frametitle{\fontsize{11}{13}\selectfont Q8: Reputational Consequences of ML}
		\begin{block}{\fontsize{8}{10}\selectfont Topic: Reputational Consequences}
			\scriptsize
			What is a reputational consequence of money laundering for a financial institution?
		\end{block}
		\vspace{0.05cm}
		\stepbox{\scriptsize
			\keypoint{Correct Answer:} Loss of customer trust, reduced business opportunities, and damage to brand value.
		}
		\vspace{0.05cm}
		\scriptsize
		\hardconcept{Core Concept:} Reputational impacts:
		\par \textbf{Trust Erosion}: Customers lose confidence.
		\par \textbf{Business Loss}: Clients and partners withdraw.
		\par \textbf{Brand Damage}: Negative media coverage.
		\par \textbf{Regulatory Scrutiny}: Increased supervision.
		\vspace{0.03cm}
		\tipbox{\scriptsize \textbf{Key Insight:} Reputational damage can be \textbf{more severe} than financial penalties.}
	\end{frame}
	
	\begin{frame}
		\frametitle{\fontsize{11}{13}\selectfont Q9: Social Consequences of ML}
		\begin{block}{\fontsize{8}{10}\selectfont Topic: Social Consequences}
			\scriptsize
			What is a social consequence of money laundering?
		\end{block}
		\vspace{0.05cm}
		\stepbox{\scriptsize
			\keypoint{Correct Answer:} Increased crime, corruption, and social inequality as illicit funds fuel criminal enterprises.
		}
		\vspace{0.05cm}
		\scriptsize
		\hardconcept{Core Concept:} Social impacts:
		\par \textbf{Crime Increase}: Laundered funds fuel more crime.
		\par \textbf{Corruption}: Criminal influence over institutions.
		\par \textbf{Inequality}: Illicit wealth concentration.
		\par \textbf{Weak Rule of Law}: Undermines legal institutions.
		\vspace{0.03cm}
		\tipbox{\scriptsize \textbf{Key Insight:} Money laundering is not a victimless crime — it fuels \textbf{organized crime} and \textbf{corruption}.}
	\end{frame}
	
	\begin{frame}
		\frametitle{\fontsize{11}{13}\selectfont Q10: Regulatory Penalties for AML Violations}
		\begin{block}{\fontsize{8}{10}\selectfont Topic: Regulatory Penalties}
			\scriptsize
			What is a potential penalty for AML/CFT regulatory violations?
		\end{block}
		\vspace{0.05cm}
		\stepbox{\scriptsize
			\keypoint{Correct Answer:} Significant financial fines, criminal prosecution of individuals, and revocation of licenses.
		}
		\vspace{0.05cm}
		\scriptsize
		\hardconcept{Core Concept:} Types of penalties:
		\par \textbf{Civil Penalties}: Financial fines (e.g., US BSA, UK FCA).
		\par \textbf{Criminal Penalties}: Imprisonment for individuals.
		\par \textbf{Administrative Penalties}: License suspension or revocation.
		\par \textbf{Personal Liability}: Compliance officers and senior management.
		\vspace{0.03cm}
		\tipbox{\scriptsize \textbf{Key Insight:} Under \textbf{FATF Recommendation 35}, sanctions must be effective, proportionate, and dissuasive.}
	\end{frame}
	
	% ================================================================
	% SLIDE 13-16: PREDICATE CRIMES
	% ================================================================
	\begin{frame}
		\frametitle{\fontsize{11}{13}\selectfont Q11: Predicate Offenses}
		\begin{block}{\fontsize{8}{10}\selectfont Topic: Predicate Crimes}
			\scriptsize
			What is a predicate offense in the context of money laundering?
		\end{block}
		\vspace{0.05cm}
		\stepbox{\scriptsize
			\keypoint{Correct Answer:} A criminal offense that generates proceeds that are then laundered through the financial system.
		}
		\vspace{0.05cm}
		\scriptsize
		\hardconcept{Core Concept:} \textbf{FATF Recommendations} define predicate offenses:
		\par \textbf{UN Convention Categories}: Drug trafficking, corruption, fraud, human trafficking.
		\par \textbf{FATF Expansion}: Tax crimes, environmental crimes, counterfeiting.
		\par \textbf{EU 6AMLD}: Expanded to include 22 predicate offenses.
		\vspace{0.03cm}
		\tipbox{\scriptsize \textbf{Key Insight:} Under FATF standards, \textbf{all serious offenses} should be predicate offenses for ML.}
	\end{frame}
	
	\begin{frame}
		\frametitle{\fontsize{11}{13}\selectfont Q12: Drug Trafficking and ML}
		\begin{block}{\fontsize{8}{10}\selectfont Topic: Drug Trafficking}
			\scriptsize
			Why is drug trafficking a major source of money laundering proceeds?
		\end{block}
		\vspace{0.05cm}
		\stepbox{\scriptsize
			\keypoint{Correct Answer:} Drug trafficking generates enormous cash proceeds that must be laundered to be used in the legitimate economy.
		}
		\vspace{0.05cm}
		\tipbox{\scriptsize \textbf{Key Insight:} The \textbf{cash-intensive nature} of drug trafficking makes it a major ML source.}
	\end{frame}
	
	\begin{frame}
		\frametitle{\fontsize{11}{13}\selectfont Q13: Corruption and ML}
		\begin{block}{\fontsize{8}{10}\selectfont Topic: Corruption}
			\scriptsize
			How does corruption relate to money laundering?
		\end{block}
		\vspace{0.05cm}
		\stepbox{\scriptsize
			\keypoint{Correct Answer:} Corruption proceeds (bribes, embezzlement) must be laundered through the financial system to appear legitimate.
		}
		\vspace{0.03cm}
		\tipbox{\scriptsize \textbf{Key Insight:} Corruption and money laundering are \textbf{interconnected} — one enables the other.}
	\end{frame}
	
	\begin{frame}
		\frametitle{\fontsize{11}{13}\selectfont Q14: Human Trafficking and ML}
		\begin{block}{\fontsize{8}{10}\selectfont Topic: Human Trafficking}
			\scriptsize
			Why is human trafficking associated with money laundering?
		\end{block}
		\vspace{0.05cm}
		\stepbox{\scriptsize
			\keypoint{Correct Answer:} Human trafficking generates significant proceeds that must be laundered through the financial system.
		}
		\vspace{0.03cm}
		\tipbox{\scriptsize \textbf{Key Insight:} Human trafficking is a high-profit, low-risk crime for traffickers, generating substantial proceeds.}
	\end{frame}
	
	% ================================================================
	% SLIDE 17-22: MANIFESTATION IN SECTORS
	% ================================================================
	\begin{frame}
		\frametitle{\fontsize{11}{13}\selectfont Q15: ML in Banking Sector}
		\begin{block}{\fontsize{8}{10}\selectfont Topic: Manifestation in Banking}
			\scriptsize
			How does money laundering manifest in the banking sector?
		\end{block}
		\vspace{0.05cm}
		\stepbox{\scriptsize
			\keypoint{Correct Answer:} Through deposits of illicit funds, layering via wire transfers, and integration through investments.
		}
		\vspace{0.05cm}
		\scriptsize
		\hardconcept{Core Concept:} Banking ML methods:
		\par \textbf{Placement}: Cash deposits, currency exchange.
		\par \textbf{Layering}: Wire transfers, multiple accounts, shell companies.
		\par \textbf{Integration}: Investments, loans, asset purchases.
		\vspace{0.03cm}
		\tipbox{\scriptsize \textbf{Key Insight:} Banking is the most common sector for money laundering due to the ease of moving funds.}
	\end{frame}
	
	\begin{frame}
		\frametitle{\fontsize{11}{13}\selectfont Q16: TF in Banking Sector}
		\begin{block}{\fontsize{8}{10}\selectfont Topic: TF in Banking}
			\scriptsize
			How does terrorist financing manifest in the banking sector?
		\end{block}
		\vspace{0.05cm}
		\stepbox{\scriptsize
			\keypoint{Correct Answer:} Through small value wire transfers, charitable donations, and use of multiple accounts to avoid detection.
		}
		\vspace{0.03cm}
		\tipbox{\scriptsize \textbf{Key Insight:} TF often involves \textbf{small amounts} to avoid triggering reporting thresholds.}
	\end{frame}
	
	\begin{frame}
		\frametitle{\fontsize{11}{13}\selectfont Q17: Sanctions Evasion in Banking}
		\begin{block}{\fontsize{8}{10}\selectfont Topic: Sanctions Evasion}
			\scriptsize
			How do banks facilitate sanctions evasion?
		\end{block}
		\vspace{0.05cm}
		\stepbox{\scriptsize
			\keypoint{Correct Answer:} Through use of shell companies, third-party intermediaries, and complex payment structures to hide sanctioned parties.
		}
		\vspace{0.03cm}
		\tipbox{\scriptsize \textbf{Key Insight:} Sanctions evasion is a \textbf{strict liability} offense — intent is not required for a violation.}
	\end{frame}
	
	\begin{frame}
		\frametitle{\fontsize{11}{13}\selectfont Q18: Bribery in Banking}
		\begin{block}{\fontsize{8}{10}\selectfont Topic: Bribery in Banking}
			\scriptsize
			How does bribery manifest in the banking sector?
		\end{block}
		\vspace{0.05cm}
		\stepbox{\scriptsize
			\keypoint{Correct Answer:} Through payments to bank officials to facilitate transactions, bypass controls, or access confidential information.
		}
		\vspace{0.03cm}
		\tipbox{\scriptsize \textbf{Key Insight:} Bribery undermines \textbf{internal controls} and enables other financial crimes.}
	\end{frame}
	
	\begin{frame}
		\frametitle{\fontsize{11}{13}\selectfont Q19: Tax Evasion in Banking}
		\begin{block}{\fontsize{8}{10}\selectfont Topic: Tax Evasion}
			\scriptsize
			How does tax evasion manifest in the banking sector?
		\end{block}
		\vspace{0.05cm}
		\stepbox{\scriptsize
			\keypoint{Correct Answer:} Through offshore accounts, shell companies, and complex structures to conceal assets and income from tax authorities.
		}
		\vspace{0.03cm}
		\tipbox{\scriptsize \textbf{Key Insight:} Tax evasion is a predicate offense for ML in many jurisdictions.}
	\end{frame}
	
	\begin{frame}
		\frametitle{\fontsize{11}{13}\selectfont Q20: Fraud in Banking}
		\begin{block}{\fontsize{8}{10}\selectfont Topic: Fraud in Banking}
			\scriptsize
			How does fraud manifest in the banking sector?
		\end{block}
		\vspace{0.05cm}
		\stepbox{\scriptsize
			\keypoint{Correct Answer:} Through unauthorized transactions, identity theft, and misrepresentation in loan or credit applications.
		}
		\vspace{0.03cm}
		\tipbox{\scriptsize \textbf{Key Insight:} Fraud proceeds must be laundered, creating a \textbf{dual crime} — fraud + ML.}
	\end{frame}
	
	% ================================================================
	% SLIDE 23-26: BANKING PRODUCT RISKS
	% ================================================================
	\begin{frame}
		\frametitle{\fontsize{11}{13}\selectfont Q21: Trade Finance Risks}
		\begin{block}{\fontsize{8}{10}\selectfont Topic: Banking Product Risks — Trade Finance}
			\scriptsize
			What is a money laundering risk in trade finance?
		\end{block}
		\vspace{0.05cm}
		\stepbox{\scriptsize
			\keypoint{Correct Answer:} Over/under-invoicing to move value across borders while disguising the true purpose of the transaction.
		}
		\vspace{0.05cm}
		\scriptsize
		\hardconcept{Core Concept:} Trade-based ML:
		\par \textbf{Over-invoicing}: Moving money out of a country.
		\par \textbf{Under-invoicing}: Moving money into a country.
		\par \textbf{Misclassification}: Changing goods descriptions.
		\par \textbf{Multiple Invoicing}: Using multiple invoices for the same goods.
		\vspace{0.03cm}
		\tipbox{\scriptsize \textbf{Key Insight:} Trade-based ML is \textbf{difficult to detect} because it appears as legitimate trade.}
	\end{frame}
	
	\begin{frame}
		\frametitle{\fontsize{11}{13}\selectfont Q22: Wealth Management Risks}
		\begin{block}{\fontsize{8}{10}\selectfont Topic: Wealth Management}
			\scriptsize
			What is a money laundering risk in wealth management?
		\end{block}
		\vspace{0.05cm}
		\stepbox{\scriptsize
			\keypoint{Correct Answer:} Use of offshore structures, trusts, and complex investment vehicles to conceal ownership and source of funds.
		}
		\vspace{0.03cm}
		\tipbox{\scriptsize \textbf{Key Insight:} Wealth management is attractive to launderers due to \textbf{complexity} and \textbf{confidentiality}.}
	\end{frame}
	
	\begin{frame}
		\frametitle{\fontsize{11}{13}\selectfont Q23: Correspondent Banking Risks}
		\begin{block}{\fontsize{8}{10}\selectfont Topic: Correspondent Banking}
			\scriptsize
			What is a risk in correspondent banking?
		\end{block}
		\vspace{0.05cm}
		\stepbox{\scriptsize
			\keypoint{Correct Answer:} The respondent bank may have weaker AML controls, creating a channel for illicit funds to enter the financial system.
		}
		\vspace{0.03cm}
		\tipbox{\scriptsize \textbf{Key Insight:} Correspondent banking is a \textbf{gateway} for cross-border money laundering.}
	\end{frame}
	
	\begin{frame}
		\frametitle{\fontsize{11}{13}\selectfont Q24: Private Banking Risks}
		\begin{block}{\fontsize{8}{10}\selectfont Topic: Private Banking}
			\scriptsize
			What is a risk in private banking?
		\end{block}
		\vspace{0.05cm}
		\stepbox{\scriptsize
			\keypoint{Correct Answer:} Discretion and personalized service can be exploited to conceal illicit activity and bypass standard controls.
		}
		\vspace{0.03cm}
		\tipbox{\scriptsize \textbf{Key Insight:} Private banking high-risk features include \textbf{numbered accounts}, \textbf{relationship discretion}, and \textbf{high asset requirements}.}
	\end{frame}
	
	% ================================================================
	% SLIDE 27-30: HIGH-RISK SECTORS
	% ================================================================
	\begin{frame}
		\frametitle{\fontsize{11}{13}\selectfont Q25: Banking High-Risk Sectors — DNFBPs}
		\begin{block}{\fontsize{8}{10}\selectfont Topic: High-Risk Sectors}
			\scriptsize
			Why is banking DNFBPs (lawyers, accountants) considered high-risk?
		\end{block}
		\vspace{0.05cm}
		\stepbox{\scriptsize
			\keypoint{Correct Answer:} DNFBPs handle client money and can facilitate transactions that conceal the source of funds.
		}
		\vspace{0.03cm}
		\tipbox{\scriptsize \textbf{Key Insight:} DNFBPs are \textbf{gatekeepers} who can enable money laundering without knowing it.}
	\end{frame}
	
	\begin{frame}
		\frametitle{\fontsize{11}{13}\selectfont Q26: Banking High-Risk Sectors — Casinos}
		\begin{block}{\fontsize{8}{10}\selectfont Topic: Casinos}
			\scriptsize
			Why is banking casinos considered high-risk?
		\end{block}
		\vspace{0.05cm}
		\stepbox{\scriptsize
			\keypoint{Correct Answer:} Casinos handle large amounts of cash and can be used for placement and layering of illicit funds.
		}
		\vspace{0.03cm}
		\tipbox{\scriptsize \textbf{Key Insight:} Casinos are \textbf{cash-intensive businesses} that can commingle illicit and legitimate funds.}
	\end{frame}
	
	\begin{frame}
		\frametitle{\fontsize{11}{13}\selectfont Q27: Banking High-Risk Sectors — Real Estate}
		\begin{block}{\fontsize{8}{10}\selectfont Topic: Real Estate}
			\scriptsize
			Why is banking real estate considered high-risk?
		\end{block}
		\vspace{0.05cm}
		\stepbox{\scriptsize
			\keypoint{Correct Answer:} Real estate purchases can be used to launder large sums of money through asset purchases.
		}
		\vspace{0.03cm}
		\tipbox{\scriptsize \textbf{Key Insight:} Real estate is attractive for \textbf{integration} — converting illicit funds into legitimate assets.}
	\end{frame}
	
	\begin{frame}
		\frametitle{\fontsize{11}{13}\selectfont Q28: Banking High-Risk Sectors — Cryptocurrency}
		\begin{block}{\fontsize{8}{10}\selectfont Topic: Cryptocurrency}
			\scriptsize
			Why is banking cryptocurrency exchanges considered high-risk?
		\end{block}
		\vspace{0.05cm}
		\stepbox{\scriptsize
			\keypoint{Correct Answer:} Pseudonymous transactions and rapid value transfer make crypto exchanges attractive for money laundering.
		}
		\vspace{0.03cm}
		\tipbox{\scriptsize \textbf{Key Insight:} Cryptocurrency exchanges are \textbf{VASPs} and are subject to FATF Recommendation 16 (Travel Rule).}
	\end{frame}
	
	% ================================================================
	% SLIDE 31-34: PEP AND HIGH-RISK CUSTOMERS
	% ================================================================
	\begin{frame}
		\frametitle{\fontsize{11}{13}\selectfont Q29: PEP Definition}
		\begin{block}{\fontsize{8}{10}\selectfont Topic: PEP and High-Risk Customers}
			\scriptsize
			What is a Politically Exposed Person (PEP)?
		\end{block}
		\vspace{0.05cm}
		\stepbox{\scriptsize
			\keypoint{Correct Answer:} An individual entrusted with a prominent public function, who is at higher risk of corruption or bribery.
		}
		\vspace{0.05cm}
		\scriptsize
		\hardconcept{Core Concept:} PEP categories:
		\par \textbf{Foreign PEPs}: From other countries (highest risk).
		\par \textbf{Domestic PEPs}: From the institution's country (lower risk).
		\par \textbf{International Organization PEPs}: In international bodies.
		\par \textbf{Family and Close Associates}: Also covered.
		\vspace{0.03cm}
		\tipbox{\scriptsize \textbf{Key Insight:} Under \textbf{FATF Recommendation 12}, institutions must have EDD for PEPs.}
	\end{frame}
	
	\begin{frame}
		\frametitle{\fontsize{11}{13}\selectfont Q30: PEP Risk Factors}
		\begin{block}{\fontsize{8}{10}\selectfont Topic: PEP Risk Factors}
			\scriptsize
			What increases the ML risk for a PEP?
		\end{block}
		\vspace{0.05cm}
		\stepbox{\scriptsize
			\keypoint{Correct Answer:} Indirect ownership structures, offshore accounts, and transactions inconsistent with known income.
		}
		\vspace{0.03cm}
		\tipbox{\scriptsize \textbf{Key Insight:} PEPs using \textbf{indirect ownership} intermediaries present the greatest concern.}
	\end{frame}
	
	\begin{frame}
		\frametitle{\fontsize{11}{13}\selectfont Q31: High-Risk Customer Indicators}
		\begin{block}{\fontsize{8}{10}\selectfont Topic: High-Risk Customers}
			\scriptsize
			What are indicators of a high-risk customer?
		\end{block}
		\vspace{0.05cm}
		\stepbox{\scriptsize
			\keypoint{Correct Answer:} Complex ownership structures, high-risk jurisdictions, unexplained wealth, and unusual transaction patterns.
		}
		\vspace{0.03cm}
		\tipbox{\scriptsize \textbf{Key Insight:} \textbf{Complexity} + \textbf{Opacity} + \textbf{High-Risk Jurisdiction} = high-risk customer.}
	\end{frame}
	
	\begin{frame}
		\frametitle{\fontsize{11}{13}\selectfont Q32: Enhanced Due Diligence for PEPs}
		\begin{block}{\fontsize{8}{10}\selectfont Topic: PEP EDD}
			\scriptsize
			What is required for PEPs under FATF Recommendation 12?
		\end{block}
		\vspace{0.05cm}
		\stepbox{\scriptsize
			\keypoint{Correct Answer:} Enhanced due diligence including source of funds/wealth verification, senior management approval, and ongoing monitoring.
		}
		\vspace{0.03cm}
		\tipbox{\scriptsize \textbf{Key Insight:} EDD for PEPs is \textbf{mandatory}, not optional.}
	\end{frame}
	
	% ================================================================
	% SLIDE 33-36: OTHER FINANCIAL SECTOR RISKS
	% ================================================================
	\begin{frame}
		\frametitle{\fontsize{11}{13}\selectfont Q33: Securities Sector Risks}
		\begin{block}{\fontsize{8}{10}\selectfont Topic: Other Financial Sector Risks}
			\scriptsize
			What is a money laundering risk in the securities sector?
		\end{block}
		\vspace{0.05cm}
		\stepbox{\scriptsize
			\keypoint{Correct Answer:} Market manipulation and insider trading generate illicit profits that must be laundered.
		}
		\vspace{0.03cm}
		\tipbox{\scriptsize \textbf{Key Insight:} The securities sector can be used for both \textbf{generating} and \textbf{laundering} illicit funds.}
	\end{frame}
	
	\begin{frame}
		\frametitle{\fontsize{11}{13}\selectfont Q34: Investment Fund Risks}
		\begin{block}{\fontsize{8}{10}\selectfont Topic: Investment Funds}
			\scriptsize
			What is a risk in investment funds?
		\end{block}
		\vspace{0.05cm}
		\stepbox{\scriptsize
			\keypoint{Correct Answer:} Use of funds to commingle illicit and legitimate investments, obscuring the source of funds.
		}
		\vspace{0.03cm}
		\tipbox{\scriptsize \textbf{Key Insight:} Investment funds can be used for \textbf{integration} of laundered funds.}
	\end{frame}
	
	\begin{frame}
		\frametitle{\fontsize{11}{13}\selectfont Q35: Leasing and Financing Risks}
		\begin{block}{\fontsize{8}{10}\selectfont Topic: Leasing and Financing}
			\scriptsize
			What is a risk in asset leasing and financing?
		\end{block}
		\vspace{0.05cm}
		\stepbox{\scriptsize
			\keypoint{Correct Answer:} Over-valuation of assets to move value, or using lease payments to layer funds.
		}
		\vspace{0.03cm}
		\tipbox{\scriptsize \textbf{Key Insight:} Asset financing can be used for \textbf{over-invoicing} and value transfer.}
	\end{frame}
	
	\begin{frame}
		\frametitle{\fontsize{11}{13}\selectfont Q36: Pension Fund Risks}
		\begin{block}{\fontsize{8}{10}\selectfont Topic: Pension Funds}
			\scriptsize
			What is a risk in pension funds?
		\end{block}
		\vspace{0.05cm}
		\stepbox{\scriptsize
			\keypoint{Correct Answer:} Illicit funds can be invested in pension funds, later withdrawn as "clean" retirement benefits.
		}
		\vspace{0.03cm}
		\tipbox{\scriptsize \textbf{Key Insight:} Pension funds provide a \textbf{long-term} integration vehicle.}
	\end{frame}
	
	% ================================================================
	% SLIDE 37-40: MSB, PSP, E-COMMERCE
	% ================================================================
	\begin{frame}
		\frametitle{\fontsize{11}{13}\selectfont Q37: MSB Risks}
		\begin{block}{\fontsize{8}{10}\selectfont Topic: MSB, PSP, E-Commerce}
			\scriptsize
			What is a money laundering risk for Money Services Businesses (MSBs)?
		\end{block}
		\vspace{0.05cm}
		\stepbox{\scriptsize
			\keypoint{Correct Answer:} Agent networks in high-risk jurisdictions with weak AML controls.
		}
		\vspace{0.03cm}
		\tipbox{\scriptsize \textbf{Key Insight:} MSB agent networks influence both \textbf{operational} and \textbf{geographic} risk exposure.}
	\end{frame}
	
	\begin{frame}
		\frametitle{\fontsize{11}{13}\selectfont Q38: PSP Risks}
		\begin{block}{\fontsize{8}{10}\selectfont Topic: PSP Risks}
			\scriptsize
			What is a risk for Payment Service Providers (PSPs)?
		\end{block}
		\vspace{0.05cm}
		\stepbox{\scriptsize
			\keypoint{Correct Answer:} Facilitating fraudulent transactions and operating in jurisdictions with limited regulatory oversight.
		}
		\vspace{0.03cm}
		\tipbox{\scriptsize \textbf{Key Insight:} PSPs are \textbf{gateways} for e-commerce fraud and money laundering.}
	\end{frame}
	
	\begin{frame}
		\frametitle{\fontsize{11}{13}\selectfont Q39: E-Commerce Platform Risks}
		\begin{block}{\fontsize{8}{10}\selectfont Topic: E-Commerce}
			\scriptsize
			What is a money laundering risk on e-commerce platforms?
		\end{block}
		\vspace{0.05cm}
		\stepbox{\scriptsize
			\keypoint{Correct Answer:} Using fake transactions or over/under-invoicing to move value across borders.
		}
		\vspace{0.03cm}
		\tipbox{\scriptsize \textbf{Key Insight:} E-commerce platforms can be used for \textbf{trade-based money laundering}.}
	\end{frame}
	
	\begin{frame}
		\frametitle{\fontsize{11}{13}\selectfont Q40: Mitigating E-Commerce Risks}
		\begin{block}{\fontsize{8}{10}\selectfont Topic: E-Commerce Mitigation}
			\scriptsize
			How can e-commerce platforms mitigate ML risk?
		\end{block}
		\vspace{0.05cm}
		\stepbox{\scriptsize
			\keypoint{Correct Answer:} Tiered onboarding based on sales activity, sanctions screening, and transaction monitoring.
		}
		\vspace{0.03cm}
		\tipbox{\scriptsize \textbf{Key Insight:} \textbf{Tiered onboarding} based on risk is an effective control.}
	\end{frame}
	
	% ================================================================
	% SLIDE 41-44: INSURANCE PRODUCT RISKS
	% ================================================================
	\begin{frame}
		\frametitle{\fontsize{11}{13}\selectfont Q41: Life Insurance Risks}
		\begin{block}{\fontsize{8}{10}\selectfont Topic: Insurance Product Risks}
			\scriptsize
			What is a money laundering risk in life insurance?
		\end{block}
		\vspace{0.05cm}
		\stepbox{\scriptsize
			\keypoint{Correct Answer:} Overpayment of premiums then early surrender, or using policies to conceal wealth.
		}
		\vspace{0.03cm}
		\tipbox{\scriptsize \textbf{Key Insight:} Life insurance can be used for \textbf{placement} and \textbf{integration}.}
	\end{frame}
	
	\begin{frame}
		\frametitle{\fontsize{11}{13}\selectfont Q42: Property/Casualty Insurance Risks}
		\begin{block}{\fontsize{8}{10}\selectfont Topic: Property/Casualty}
			\scriptsize
			What is a risk in property/casualty insurance?
		\end{block}
		\vspace{0.05cm}
		\stepbox{\scriptsize
			\keypoint{Correct Answer:} Fraudulent claims to convert illicit funds into legitimate claim payments.
		}
		\vspace{0.03cm}
		\tipbox{\scriptsize \textbf{Key Insight:} Insurance fraud generates proceeds that need laundering.}
	\end{frame}
	
	\begin{frame}
		\frametitle{\fontsize{11}{13}\selectfont Q43: Reinsurance Risks}
		\begin{block}{\fontsize{8}{10}\selectfont Topic: Reinsurance}
			\scriptsize
			What is a risk in reinsurance?
		\end{block}
		\vspace{0.05cm}
		\stepbox{\scriptsize
			\keypoint{Correct Answer:} Complex transactions that obscure the flow of funds across borders.
		}
		\vspace{0.03cm}
		\tipbox{\scriptsize \textbf{Key Insight:} Reinsurance complexity can hide \textbf{cross-border} fund movement.}
	\end{frame}
	
	\begin{frame}
		\frametitle{\fontsize{11}{13}\selectfont Q44: Insurance CDD Requirements}
		\begin{block}{\fontsize{8}{10}\selectfont Topic: Insurance CDD}
			\scriptsize
			What CDD is required for insurance products under FATF?
		\end{block}
		\vspace{0.05cm}
		\stepbox{\scriptsize
			\keypoint{Correct Answer:} Identification of the policyholder and beneficiary, with EDD for high-risk policies.
		}
		\vspace{0.03cm}
		\tipbox{\scriptsize \textbf{Key Insight:} Under \textbf{FATF Recommendation 10}, insurance companies are covered entities.}
	\end{frame}
	
	% ================================================================
	% SLIDE 45-48: VASP AND CRYPTOASSETS
	% ================================================================
	\begin{frame}
		\frametitle{\fontsize{11}{13}\selectfont Q45: VASP Definition}
		\begin{block}{\fontsize{8}{10}\selectfont Topic: VASP and Cryptoassets}
			\scriptsize
			What is a Virtual Asset Service Provider (VASP)?
		\end{block}
		\vspace{0.05cm}
		\stepbox{\scriptsize
			\keypoint{Correct Answer:} A person or entity conducting virtual asset exchange, transfer, or custody services for third parties.
		}
		\vspace{0.03cm}
		\tipbox{\scriptsize \textbf{Key Insight:} FATF extends AML requirements to VASPs under \textbf{Recommendation 16}.}
	\end{frame}
	
	\begin{frame}
		\frametitle{\fontsize{11}{13}\selectfont Q46: Cryptoasset ML Risks}
		\begin{block}{\fontsize{8}{10}\selectfont Topic: Cryptoasset Risks}
			\scriptsize
			What is a money laundering risk in cryptoassets?
		\end{block}
		\vspace{0.05cm}
		\stepbox{\scriptsize
			\keypoint{Correct Answer:} Pseudonymity, rapid cross-border transfer, and use of privacy coins/mixers to obscure transactions.
		}
		\vspace{0.03cm}
		\tipbox{\scriptsize \textbf{Key Insight:} Privacy coins (Monero, Zcash) and mixers (Tornado Cash) are \textbf{challenges} for tracing.}
	\end{frame}
	
	\begin{frame}
		\frametitle{\fontsize{11}{13}\selectfont Q47: Travel Rule for VASPs}
		\begin{block}{\fontsize{8}{10}\selectfont Topic: Travel Rule}
			\scriptsize
			What is the Travel Rule requirement for VASPs?
		\end{block}
		\vspace{0.05cm}
		\stepbox{\scriptsize
			\keypoint{Correct Answer:} VASPs must collect and share originator and beneficiary information for virtual asset transfers.
		}
		\vspace{0.03cm}
		\tipbox{\scriptsize \textbf{Key Insight:} The Travel Rule applies to VASPs under \textbf{FATF Recommendation 16}.}
	\end{frame}
	
	\begin{frame}
		\frametitle{\fontsize{11}{13}\selectfont Q48: DeFi and ML Risks}
		\begin{block}{\fontsize{8}{10}\selectfont Topic: DeFi Risks}
			\scriptsize
			What is a risk in Decentralized Finance (DeFi)?
		\end{block}
		\vspace{0.05cm}
		\stepbox{\scriptsize
			\keypoint{Correct Answer:} Lack of centralized oversight makes tracing and AML controls difficult.
		}
		\vspace{0.03cm}
		\tipbox{\scriptsize \textbf{Key Insight:} DeFi platforms can be used for \textbf{rapid, anonymous} value transfer.}
	\end{frame}
	
	% ================================================================
	% SLIDE 49-52: GAMING AND GAMBLING
	% ================================================================
	\begin{frame}
		\frametitle{\fontsize{11}{13}\selectfont Q49: Casino ML Risks}
		\begin{block}{\fontsize{8}{10}\selectfont Topic: Gaming and Gambling}
			\scriptsize
			What is a money laundering risk in casinos?
		\end{block}
		\vspace{0.05cm}
		\stepbox{\scriptsize
			\keypoint{Correct Answer:} Cash purchase of chips, minimal play, then cash out with clean funds (placement/layering).
		}
		\vspace{0.03cm}
		\tipbox{\scriptsize \textbf{Key Insight:} Casinos can be used for both \textbf{placement} and \textbf{layering}.}
	\end{frame}
	
	\begin{frame}
		\frametitle{\fontsize{11}{13}\selectfont Q50: Online Gambling Risks}
		\begin{block}{\fontsize{8}{10}\selectfont Topic: Online Gambling}
			\scriptsize
			What is a risk in online gambling?
		\end{block}
		\vspace{0.05cm}
		\stepbox{\scriptsize
			\keypoint{Correct Answer:} Deposits and withdrawals via multiple payment methods, obscuring the source of funds.
		}
		\vspace{0.03cm}
		\tipbox{\scriptsize \textbf{Key Insight:} Online gambling is \textbf{high-risk} due to cross-border payment complexity.}
	\end{frame}
	
	\begin{frame}
		\frametitle{\fontsize{11}{13}\selectfont Q51: Sports Betting Risks}
		\begin{block}{\fontsize{8}{10}\selectfont Topic: Sports Betting}
			\scriptsize
			What is a risk in sports betting?
		\end{block}
		\vspace{0.05cm}
		\stepbox{\scriptsize
			\keypoint{Correct Answer:} Match-fixing and insider betting generate proceeds that need laundering.
		}
		\vspace{0.03cm}
		\tipbox{\scriptsize \textbf{Key Insight:} Sports betting can be used to \textbf{generate} and \textbf{launder} funds.}
	\end{frame}
	
	\begin{frame}
		\frametitle{\fontsize{11}{13}\selectfont Q52: Gaming CDD Requirements}
		\begin{block}{\fontsize{8}{10}\selectfont Topic: Gaming CDD}
			\scriptsize
			What CDD is required for gaming establishments?
		\end{block}
		\vspace{0.05cm}
		\stepbox{\scriptsize
			\keypoint{Correct Answer:} Identification at set thresholds, transaction monitoring, and STR filing.
		}
		\vspace{0.03cm}
		\tipbox{\scriptsize \textbf{Key Insight:} Gaming establishments are \textbf{covered entities} under FATF.}
	\end{frame}
	
	% ================================================================
	% SLIDE 53-56: REAL ESTATE
	% ================================================================
	\begin{frame}
		\frametitle{\fontsize{11}{13}\selectfont Q53: Real Estate ML Risks}
		\begin{block}{\fontsize{8}{10}\selectfont Topic: Real Estate}
			\scriptsize
			What is a money laundering risk in real estate?
		\end{block}
		\vspace{0.05cm}
		\stepbox{\scriptsize
			\keypoint{Correct Answer:} Purchase of property with illicit funds, over/under-valuation, and use of shell companies to hide ownership.
		}
		\vspace{0.03cm}
		\tipbox{\scriptsize \textbf{Key Insight:} Real estate is a preferred \textbf{integration} vehicle.}
	\end{frame}
	
	\begin{frame}
		\frametitle{\fontsize{11}{13}\selectfont Q54: Real Estate Shell Companies}
		\begin{block}{\fontsize{8}{10}\selectfont Topic: Shell Companies in Real Estate}
			\scriptsize
			Why are shell companies used in real estate money laundering?
		\end{block}
		\vspace{0.05cm}
		\stepbox{\scriptsize
			\keypoint{Correct Answer:} They conceal beneficial ownership and make tracing the source of funds difficult.
		}
		\vspace{0.03cm}
		\tipbox{\scriptsize \textbf{Key Insight:} Shell companies in real estate are a \textbf{key red flag}.}
	\end{frame}
	
	\begin{frame}
		\frametitle{\fontsize{11}{13}\selectfont Q55: Real Estate Over-Valuation}
		\begin{block}{\fontsize{8}{10}\selectfont Topic: Over-Valuation}
			\scriptsize
			How is over-valuation used in real estate money laundering?
		\end{block}
		\vspace{0.05cm}
		\stepbox{\scriptsize
			\keypoint{Correct Answer:} Property is purchased above market value, with the difference returned to the buyer as "clean" proceeds.
		}
		\vspace{0.03cm}
		\tipbox{\scriptsize \textbf{Key Insight:} Over-valuation is a form of \textbf{value transfer}.}
	\end{frame}
	
	\begin{frame}
		\frametitle{\fontsize{11}{13}\selectfont Q56: Real Estate EDD Requirements}
		\begin{block}{\fontsize{8}{10}\selectfont Topic: Real Estate EDD}
			\scriptsize
			What EDD is required for high-risk real estate transactions?
		\end{block}
		\vspace{0.05cm}
		\stepbox{\scriptsize
			\keypoint{Correct Answer:} Source of funds verification, beneficial ownership identification, and transaction monitoring.
		}
		\vspace{0.03cm}
		\tipbox{\scriptsize \textbf{Key Insight:} Real estate agents are covered under \textbf{FATF Recommendation 22}.}
	\end{frame}
	
	% ================================================================
	% SLIDE 57-60: GATEKEEPER RISKS
	% ================================================================
	\begin{frame}
		\frametitle{\fontsize{11}{13}\selectfont Q57: Gatekeeper Definition}
		\begin{block}{\fontsize{8}{10}\selectfont Topic: Gatekeeper Risks}
			\scriptsize
			What is a gatekeeper in the AML context?
		\end{block}
		\vspace{0.05cm}
		\stepbox{\scriptsize
			\keypoint{Correct Answer:} Professionals who facilitate financial transactions or provide services that can be exploited for money laundering.
		}
		\vspace{0.03cm}
		\tipbox{\scriptsize \textbf{Key Insight:} Gatekeepers include lawyers, accountants, and TCSPs.}
	\end{frame}
	
	\begin{frame}
		\frametitle{\fontsize{11}{13}\selectfont Q58: Lawyer Gatekeeper Risks}
		\begin{block}{\fontsize{8}{10}\selectfont Topic: Lawyer Risks}
			\scriptsize
			What is a risk with lawyers as gatekeepers?
		\end{block}
		\vspace{0.05cm}
		\stepbox{\scriptsize
			\keypoint{Correct Answer:} Handling client funds in escrow accounts, creating structures to conceal ownership.
		}
		\vspace{0.03cm}
		\tipbox{\scriptsize \textbf{Key Insight:} Lawyers are covered under \textbf{FATF Recommendation 22}.}
	\end{frame}
	
	\begin{frame}
		\frametitle{\fontsize{11}{13}\selectfont Q59: Notary Gatekeeper Risks}
		\begin{block}{\fontsize{8}{10}\selectfont Topic: Notary Risks}
			\scriptsize
			What is a risk with notaries as gatekeepers?
		\end{block}
		\vspace{0.05cm}
		\stepbox{\scriptsize
			\keypoint{Correct Answer:} Authenticating documents that may be used to conceal or transfer ownership of illicit assets.
		}
		\vspace{0.03cm}
		\tipbox{\scriptsize \textbf{Key Insight:} Notaries are covered under \textbf{FATF Recommendation 22}.}
	\end{frame}
	
	\begin{frame}
		\frametitle{\fontsize{11}{13}\selectfont Q60: Mitigating Gatekeeper Risks}
		\begin{block}{\fontsize{8}{10}\selectfont Topic: Mitigating Gatekeeper Risks}
			\scriptsize
			How can gatekeeper risks be mitigated?
		\end{block}
		\vspace{0.05cm}
		\stepbox{\scriptsize
			\keypoint{Correct Answer:} Training, client acceptance committees, and standardized due diligence questionnaires.
		}
		\vspace{0.03cm}
		\tipbox{\scriptsize \textbf{Key Insight:} \textbf{Training} and \textbf{client acceptance committees} are effective controls.}
	\end{frame}
	
	% ================================================================
	% SLIDE 61-64: TCSP RISKS
	% ================================================================
	\begin{frame}
		\frametitle{\fontsize{11}{13}\selectfont Q61: TCSP Definition}
		\begin{block}{\fontsize{8}{10}\selectfont Topic: Trust and Company Service Providers}
			\scriptsize
			What is a Trust and Company Service Provider (TCSP)?
		\end{block}
		\vspace{0.05cm}
		\stepbox{\scriptsize
			\keypoint{Correct Answer:} An entity providing services such as company formation, registered office, nominee directors, and trust services.
		}
		\vspace{0.03cm}
		\tipbox{\scriptsize \textbf{Key Insight:} TCSPs are covered under \textbf{FATF Recommendation 23}.}
	\end{frame}
	
	\begin{frame}
		\frametitle{\fontsize{11}{13}\selectfont Q62: TCSP ML Risks}
		\begin{block}{\fontsize{8}{10}\selectfont Topic: TCSP Risks}
			\scriptsize
			What is a money laundering risk in TCSP services?
		\end{block}
		\vspace{0.05cm}
		\stepbox{\scriptsize
			\keypoint{Correct Answer:} Formation of companies to conceal beneficial ownership and facilitate layering of illicit funds.
		}
		\vspace{0.03cm}
		\tipbox{\scriptsize \textbf{Key Insight:} TCSPs can create \textbf{vehicles} for money laundering.}
	\end{frame}
	
	\begin{frame}
		\frametitle{\fontsize{11}{13}\selectfont Q63: TCSP Red Flags}
		\begin{block}{\fontsize{8}{10}\selectfont Topic: TCSP Red Flags}
			\scriptsize
			What are red flags for TCSP clients?
		\end{block}
		\vspace{0.05cm}
		\stepbox{\scriptsize
			\keypoint{Correct Answer:} Requesting confidentiality for shareholders, high-risk jurisdictions, and refusal to disclose beneficial ownership.
		}
		\vspace{0.03cm}
		\tipbox{\scriptsize \textbf{Key Insight:} TCSPs should \textbf{decline} clients who refuse UBO disclosure.}
	\end{frame}
	
	\begin{frame}
		\frametitle{\fontsize{11}{13}\selectfont Q64: TCSP CDD Requirements}
		\begin{block}{\fontsize{8}{10}\selectfont Topic: TCSP CDD}
			\scriptsize
			What CDD is required for TCSPs?
		\end{block}
		\vspace{0.05cm}
		\stepbox{\scriptsize
			\keypoint{Correct Answer:} Identify the client and beneficial owner, understand the purpose, and monitor the relationship.
		}
		\vspace{0.03cm}
		\tipbox{\scriptsize \textbf{Key Insight:} CDD for TCSPs is \textbf{mandatory} under FATF.}
	\end{frame}
	
	% ================================================================
	% SLIDE 65-68: ACCOUNTANT RISKS
	% ================================================================
	\begin{frame}
		\frametitle{\fontsize{11}{13}\selectfont Q65: Accountant Gatekeeper Risks}
		\begin{block}{\fontsize{8}{10}\selectfont Topic: Accountant Risks}
			\scriptsize
			What is a money laundering risk with accountants?
		\end{block}
		\vspace{0.05cm}
		\stepbox{\scriptsize
			\keypoint{Correct Answer:} Preparing accounts that conceal the source of funds and facilitating tax evasion.
		}
		\vspace{0.03cm}
		\tipbox{\scriptsize \textbf{Key Insight:} Accountants can be used to \textbf{conceal} illicit activity.}
	\end{frame}
	
	\begin{frame}
		\frametitle{\fontsize{11}{13}\selectfont Q66: Accountant Red Flags}
		\begin{block}{\fontsize{8}{10}\selectfont Topic: Accountant Red Flags}
			\scriptsize
			What are red flags for accountants?
		\end{block}
		\vspace{0.05cm}
		\stepbox{\scriptsize
			\keypoint{Correct Answer:} Clients who refuse to provide source of funds, complex structures without clear business purpose.
		}
		\vspace{0.03cm}
		\tipbox{\scriptsize \textbf{Key Insight:} Accountants must \textbf{question} unusual client requests.}
	\end{frame}
	
	\begin{frame}
		\frametitle{\fontsize{11}{13}\selectfont Q67: Accountant ML Facilitation}
		\begin{block}{\fontsize{8}{10}\selectfont Topic: Accountant Facilitation}
			\scriptsize
			How can accountants facilitate money laundering?
		\end{block}
		\vspace{0.05cm}
		\stepbox{\scriptsize
			\keypoint{Correct Answer:} Creating false records, overstating expenses, and concealing beneficial ownership.
		}
		\vspace{0.03cm}
		\tipbox{\scriptsize \textbf{Key Insight:} Accountants can be \textbf{willing} or \textbf{unwilling} facilitators.}
	\end{frame}
	
	\begin{frame}
		\frametitle{\fontsize{11}{13}\selectfont Q68: Accountant Responsibilities}
		\begin{block}{\fontsize{8}{10}\selectfont Topic: Accountant Responsibilities}
			\scriptsize
			What are accountants' AML responsibilities?
		\end{block}
		\vspace{0.05cm}
		\stepbox{\scriptsize
			\keypoint{Correct Answer:} CDD, record-keeping, and STR filing for suspicious activity.
		}
		\vspace{0.03cm}
		\tipbox{\scriptsize \textbf{Key Insight:} Accountants are covered under \textbf{FATF Recommendation 22}.}
	\end{frame}
	
	% ================================================================
	% SLIDE 69-72: OTHER HIGH-RISK SECTORS
	% ================================================================
	\begin{frame}
		\frametitle{\fontsize{11}{13}\selectfont Q69: Free-Trade Zone Risks}
		\begin{block}{\fontsize{8}{10}\selectfont Topic: Other High-Risk Sectors}
			\scriptsize
			What is a money laundering risk in Free-Trade Zones (FTZs)?
		\end{block}
		\vspace{0.05cm}
		\stepbox{\scriptsize
			\keypoint{Correct Answer:} Lack of standardized inspection protocols for transshipped goods, enabling trade-based ML.
		}
		\vspace{0.03cm}
		\tipbox{\scriptsize \textbf{Key Insight:} FTZs are \textbf{high-risk} due to limited oversight.}
	\end{frame}
	
	\begin{frame}
		\frametitle{\fontsize{11}{13}\selectfont Q70: Precious Metals and Stones Risks}
		\begin{block}{\fontsize{8}{10}\selectfont Topic: Precious Metals/Stones}
			\scriptsize
			What is a risk in precious metals and stones?
		\end{block}
		\vspace{0.05cm}
		\stepbox{\scriptsize
			\keypoint{Correct Answer:} High-value, portable, and easily concealed assets that can be used for value transfer.
		}
		\vspace{0.03cm}
		\tipbox{\scriptsize \textbf{Key Insight:} Precious metals are covered under \textbf{FATF Recommendation 23}.}
	\end{frame}
	
	\begin{frame}
		\frametitle{\fontsize{11}{13}\selectfont Q71: Art and Antiquities Risks}
		\begin{block}{\fontsize{8}{10}\selectfont Topic: Art and Antiquities}
			\scriptsize
			What is a risk in art and antiquities?
		\end{block}
		\vspace{0.05cm}
		\stepbox{\scriptsize
			\keypoint{Correct Answer:} Subjective valuation makes over/under-invoicing easy, and sales can be anonymous.
		}
		\vspace{0.03cm}
		\tipbox{\scriptsize \textbf{Key Insight:} Art is attractive for \textbf{value transfer} and \textbf{integration}.}
	\end{frame}
	
	\begin{frame}
		\frametitle{\fontsize{11}{13}\selectfont Q72: Shell Companies and Trusts}
		\begin{block}{\fontsize{8}{10}\selectfont Topic: Shell Companies/Trusts}
			\scriptsize
			What is the risk of shell companies and trusts?
		\end{block}
		\vspace{0.05cm}
		\stepbox{\scriptsize
			\keypoint{Correct Answer:} They conceal beneficial ownership and create opacity in financial transactions.
		}
		\vspace{0.03cm}
		\tipbox{\scriptsize \textbf{Key Insight:} Shell companies are a \textbf{layering} tool.}
	\end{frame}
	
	% ================================================================
	% SLIDE 73-80: SUMMARY AND THANK YOU
	% ================================================================
	\begin{frame}
		\frametitle{\fontsize{11}{13}\selectfont SUMMARY: Key Principles — Domain A}
		\scriptsize
		\begin{block}{\fontsize{9}{11}\selectfont Understanding the Risks and Methods of Financial Crime:}
			\vspace{0.05cm}
			\scriptsize
			\begin{tabular}{|p{0.35\textwidth}|p{0.60\textwidth}|}
				\hline
				\keypoint{Concept} & \keypoint{Application} \\
				\hline
				Money Laundering & 3 stages: Placement, Layering, Integration \\
				\hline
				Terrorist Financing & Funds may be legitimate; intent is key \\
				\hline
				Sanctions & Strict liability; OFAC, UN, EU regimes \\
				\hline
				Fraud & Identity theft, synthetic fraud, payment fraud \\
				\hline
				Anti-Bribery & UK Act prohibits facilitation; FCPA allows \\
				\hline
				Tax Evasion & Illegal concealment; predicate for ML \\
				\hline
				ML Consequences & Economic, reputational, social \\
				\hline
				Predicate Crimes & Drug trafficking, corruption, human trafficking \\
				\hline
				Banking Products & Trade finance, wealth, correspondent, private \\
				\hline
				High-Risk Sectors & DNFBPs, casinos, real estate, crypto \\
				\hline
				PEPs & Foreign, domestic, int'l org; EDD mandatory \\
				\hline
				MSB/PSP/E-Commerce & Agent networks, fraud facilitation \\
				\hline
				Insurance & Life, property/casualty, reinsurance risks \\
				\hline
				VASP/Crypto & Travel Rule, privacy coins, DeFi risks \\
				\hline
				Gaming/Gambling & Casinos, online, sports betting \\
				\hline
				Real Estate & Shell companies, over-valuation \\
				\hline
				Gatekeepers & Lawyers, accountants, TCSPs \\
				\hline
				TCSPs & Company formation, nominee directors \\
				\hline
				FTZs & Limited inspection; trade-based ML \\
				\hline
				Precious Metals/Art & Subjective valuation; value transfer \\
				\hline
			\end{tabular}
			\vspace{0.05cm}
		\end{block}
		\vspace{0.05cm}
		\tipbox{\scriptsize \textbf{FINAL LESSON:} Domain A tests understanding of \textbf{how financial crime works} — the definitions, methods, and sector-specific risks. Understand the \textit{methods} and the \textit{answers} become obvious.}
	\end{frame}
	
	\begin{frame}
		\frametitle{\fontsize{11}{13}\selectfont EXAM TIPS — Domain A}
		\scriptsize
		\begin{block}{\fontsize{9}{11}\selectfont Common Traps to Avoid:}
			\vspace{0.05cm}
			\scriptsize
			\begin{itemize}
				\item \highlight{TF vs ML}: TF funds can be \textbf{legitimate}; ML funds are \textbf{always illicit}.
				\item \highlight{Tax Evasion vs Avoidance}: Evasion is \textbf{illegal}; avoidance is \textbf{legal}.
				\item \highlight{Facilitation Payments}: Allowed under FCPA; \textbf{prohibited} under UK Bribery Act.
				\item \highlight{Placement vs Layering}: Placement = \textbf{entering} the system; Layering = \textbf{moving} to obscure origin.
				\item \highlight{Integration}: Converting illicit funds into \textbf{legitimate} assets.
				\item \highlight{PEP EDD}: EDD is \textbf{mandatory} under FATF Recommendation 12.
				\item \highlight{VASP Travel Rule}: VASPs must share originator/beneficiary information.
				\item \highlight{FTZ Risk}: \textbf{Lack of inspection} is the key vulnerability.
				\item \highlight{Gatekeepers}: Lawyers, accountants, TCSPs are \textbf{covered entities}.
			\end{itemize}
			\vspace{0.05cm}
		\end{block}
		\vspace{0.05cm}
		\tipbox{\scriptsize \textbf{Key Insight:} Focus on understanding the \textit{methods} criminals use in each sector.}
	\end{frame}
	
	\begin{frame}
		\centering
		\vspace{2.5cm}
		{\fontsize{16}{19}\selectfont \textbf{Thank You}} \\[0.2cm]
		{\fontsize{11}{13}\selectfont Keep learning. Keep growing.} \\[0.15cm]
		\rule{3cm}{0.5pt} \\[0.15cm]
		{\fontsize{8}{10}\selectfont \textit{Understanding financial crime risks is the foundation of effective AML/CFT compliance.}}
		\vspace{0.3cm}
		{\fontsize{7}{9}\selectfont \textcolor{darkgray}{End of Domain A Review — 80 Slides}}
	\end{frame}
	
\end{document}